% Auto-generated by src/analysis/latex_tables.py — do not edit by hand.
\begin{table}[htbp]
\centering
\small
\caption{Top-10 donor counties for the SDID synthetic control (headline cath\_share $>$ 50\% specification)}
\label{tab:sdid_donor_counties}
\begin{tabular}{rlcrr}
\toprule
 & Kreis & Rb & Cath.\ share (\%) & $\hat\omega$ \\
\midrule
\multicolumn{5}{l}{\textit{Panel A: Marriage-rate synthetic control}} \\
1 & Reichenbach & BRE & 29.3 & 0.0074 \\
2 & Weissensee & ERF & 1.1 & 0.0073 \\
3 & Schleusingen & ERF & 0.4 & 0.0069 \\
4 & Langensalza & ERF & 0.7 & 0.0068 \\
5 & Neidenburg & KON & 14.1 & 0.0068 \\
6 & Ziegenruck & ERF & 0.3 & 0.0067 \\
7 & Lennep & DUS & 16.9 & 0.0066 \\
8 & Wittgenstein & ARN & 3.3 & 0.0066 \\
9 & Eckartsberga & MER & 0.2 & 0.0065 \\
10 & Gummersbach & KOL & 8.2 & 0.0064 \\
\midrule\multicolumn{5}{l}{Effective number of donors $1/\sum_i \hat\omega_i^2 = 172.0$; max $\hat\omega = 0.0074$} \\
\midrule
\multicolumn{5}{l}{\textit{Panel B: CBR synthetic control}} \\
1 & Lennep & DUS & 16.9 & 0.0093 \\
2 & Franzburg & STR & 0.1 & 0.0077 \\
3 & Reichenbach & BRE & 29.3 & 0.0076 \\
4 & Muhlhausen & ERF & 28.2 & 0.0075 \\
5 & Zauch-Belzig & POT & 0.4 & 0.0074 \\
6 & Waldenburg & BRE & 22.9 & 0.0073 \\
7 & Weissensee & ERF & 1.1 & 0.0072 \\
8 & Siegen & ARN & 17.0 & 0.0072 \\
9 & Halle & MIN & 1.9 & 0.0072 \\
10 & Naumburg & MER & 0.9 & 0.0071 \\
\midrule\multicolumn{5}{l}{Effective number of donors $1/\sum_i \hat\omega_i^2 = 167.6$; max $\hat\omega = 0.0093$} \\
\bottomrule
\end{tabular}
\begin{minipage}{\linewidth}
\vspace{0.5em}
\footnotesize \textit{Notes:} Unit weights $\hat\omega_i$ are the simplex-constrained ridge least-squares solution that best matches the high-Catholic pre-period trajectory using a convex combination of low-Catholic counties. The effective number of donors $1/\sum_i \hat\omega_i^2$ summarises weight diffusion: values close to the total number of control counties indicate near-uniform weighting and rule out concentration risk (i.e.\ the result is not driven by a handful of idiosyncratic donors). Rb $=$ Regierungsbezirk (Prussian administrative region). The donor pool spans rural and mixed-agrarian Prussian Kreise across Rheinland (DUS, KOB), Westphalia (ARN, MIN), Silesia (BRE, LIE), and Thuringia/Saxony (ERF, MER) -- structurally comparable to the Catholic Kreise of Westphalia and the Rhineland, which mitigates concerns about geographic compositional bias.
\end{minipage}
\end{table}
